% CPAIOR 2027 short paper -- draft 3 (2026-10-23).
% Story: a policy learned from the look-ahead's action costs, and a learned
% decision of WHEN to call the look-ahead (gate + rollout check).
% Previous draft (2026-10-04): main_draft2.tex.  Needs llncs.cls and
% splncs04.bst from the user's template, and references.bib.
% Every number comes from ML/results; the script is named above each table.
% Red \tbd{} marks results that are still running (mixed-charger routes).
\documentclass[runningheads]{llncs}
%
\usepackage[T1]{fontenc}
\usepackage{newtxtext}
\usepackage[varvw]{newtxmath}
\usepackage{graphicx}
\usepackage{booktabs}
\usepackage{algorithm}
\usepackage{algpseudocode}
\usepackage{xcolor}
\newcommand{\tbd}[1]{\textcolor{red}{[#1]}}
%
\begin{document}
%
\title{Learning When to Call the Solver: Imitation and Selective Look-Ahead
for Charging and Rest Decisions of Electric Trucks}
\titlerunning{Learning When to Call the Solver}
%
\author{C\'eline Pagnier\inst{1}\orcidID{0009-0009-7276-9525} \and
Steffen J. S. Bakker\inst{2,3}\orcidID{0000-0002-2656-3827}}
%
\authorrunning{C. Pagnier and S. J. S. Bakker}
% TODO: affiliations 2 and 3 for the second author
\institute{Department of Industrial Economics and Technology Management,
Norwegian University of Science and Technology, Trondheim, Norway}
%
\maketitle
%
\begin{abstract}
We consider the online scheduling of charging stops, breaks and daily rests
of a battery-electric truck under the EU driving-time regulation and uncertain
travel times. A stochastic rolling-horizon look-ahead policy, which solves a
mixed-integer program for every candidate action under sampled travel-time
scenarios, performs well but requires more than a minute of solver time per
decision. We first learn a policy from the costs that the look-ahead computes
for all candidate actions and execute it behind a rule-based safety filter. On
125 held-out routes, the learned policies are within 0.7 percentage points of
the look-ahead and violate the regulation in none of 750 runs, whereas a policy
learned from
the hindsight-optimal schedules is markedly worse. We then study when calling
the look-ahead is worth its cost. A learned gate, trained on out-of-fold
errors of the learned policy, selects the decisions that are passed to the
look-ahead, and a rollout check accepts the look-ahead's action only if it
remains better when the learned policy continues the route. On the held-out
routes, this hybrid policy matches the look-ahead ($-0.36\pm0.33$ points) and
improves on the learned policy by $1.09\pm0.30$ points while calling the
look-ahead at 8.5\% of the decisions. Both components are needed: without the
rollout check, actions of the look-ahead that are followed by the learned
policy cause additional daily rests. \tbd{One sentence on charging stations of
varying power.}
\keywords{Imitation learning \and Rolling-horizon look-ahead \and Hybrid
policies \and Electric vehicles \and Driving-time regulation}
\end{abstract}
%
%
\section{Introduction}
\label{sec:intro}

The electrification of road freight requires heavy-duty trucks to recharge
during long-haul trips. For a battery-electric truck, charging cannot be
planned independently of the driver: EU Regulation 561/2006~\cite{eu561}
imposes breaks after 4.5 hours of driving and daily rests within each working
day, and these stops are natural opportunities to charge. Since travel times
are uncertain, the decisions of when to charge, for how long, and when to
take a break or a rest are taken online, at each stop of the route.

In a companion paper~\cite{pagnier_bakker_journal}, we formulated this joint
problem as a mixed-integer program (MIP) and compared several online policies.
The best of them is a stochastic rolling-horizon look-ahead (LA), which
evaluates every candidate action at the current stop by solving the MIP over
the next 24 hours for 25 sampled travel-time scenarios. The LA stays within
about 3\% of the hindsight-optimal schedule, but a decision requires more than
a minute of solver time with a commercial MIP solver, and the computation is
repeated at every stop of every route. A learned policy decides in
milliseconds, but it can fail on situations that its training data do not
cover, for instance on routes whose charging stations have different powers.
In this paper, we therefore do not only ask whether a learned policy can
replace the LA, but also when the LA is worth calling.

Our contributions are as follows. First, we learn a policy from the costs that
the LA computes for all candidate actions, and execute it behind a rule-based
safety filter that enforces the regulation independently of the learned
model. With gradient-boosted trees as well as a neural network, the learned
policy is within 0.7 percentage points of the LA on 125 held-out routes,
whereas the same models trained on the hindsight-optimal schedules are
markedly worse. Second, we propose a hybrid policy that calls the LA
selectively. A gate, trained on out-of-fold errors of the learned policy,
decides at which stops to call the LA, and a rollout check accepts the LA's
action only if it remains better when the learned policy continues the route.
On the held-out routes, the hybrid policy matches the LA while calling it at
8.5\% of the decisions, and an ablation shows that both the gate and the
rollout check contribute. Third, we evaluate the hybrid policy on routes whose
charging stations have different powers, where the learned policy alone is
$2.84\pm1.04$ percentage points behind the LA. \tbd{Result.}

Section~\ref{sec:related} reviews related work, Sect.~\ref{sec:problem}
defines the problem and the LA, Sect.~\ref{sec:method} presents the learned
and the hybrid policies, and Sect.~\ref{sec:exp} reports the computational
experiments. Section~\ref{sec:conclusion} concludes.

\section{Related Work}
\label{sec:related}

\subsubsection{Learning from optimization experts.}
Imitation learning trains a policy on states labelled by an expert, and
DAgger~\cite{ross2011} collects these states with the learned policy to
counter the drift away from the expert's state distribution.
AggreVaTe~\cite{ross_bagnell2014} and its differentiable
variant~\cite{sun2017aggrevated} use the expert's cost-to-go of every action
instead of its choice only, which bounds the performance of the learned policy
by that of the expert plus the error of the cost regression. Our training
targets are of this kind: the LA's scenario-average cost of every action. As
noted in the analysis of LOLS~\cite{chang2015lols}, such guarantees are
relative to the expert, and improving on it requires roll-outs with the
learned policy. The LA itself evaluates actions on sampled futures that are
solved with perfect information, as the expectation algorithm of online
stochastic combinatorial optimization~\cite{vanhentenryck2006} and hindsight
optimization~\cite{yoon2008}. Policies can also be learned from
hindsight-optimal solutions directly~\cite{baty2024,sinclair2023}; we compare
both teachers. At CPAIOR, Bourgeat et al.~\cite{bourgeat2026} imitate a global
solver before fine-tuning with reinforcement learning, and Anand et
al.~\cite{anand2026} fine-tune a MIP-based policy with reinforcement learning.

\subsubsection{Approximate predictive control and safety.}
The LA is a stochastic model predictive controller. Learned approximations of
predictive controllers are studied in control~\cite{hertneck2018,karg2020},
often together with a safety filter that modifies unsafe inputs of the learned
controller~\cite{wabersich2021}, or a shield that restricts the actions of a
learned policy~\cite{alshiekh2018}. Zhang et al.~\cite{zhang2019certificate}
learn, besides the policy, an estimate of its suboptimality that is evaluated
online. Deciding between a cheap and an expensive decision procedure is also
related to algorithm selection, which Guigal et al.~\cite{guigal2026} apply to
heuristics and an exact method inside a rolling horizon.

\subsubsection{Electric vehicles and driving-time regulations.}
Electric vehicle routing with recharging stations~\cite{schneider2014} and
nonlinear charging functions~\cite{montoya2017}, and vehicle scheduling under
drivers' working hours~\cite{goel2009,goel_vidal2014}, have mostly been
studied separately and with deterministic travel times. Dynamic and
stochastic variants of electric vehicle routing are addressed with rollout
policies~\cite{kullman2021} and safe reinforcement learning~\cite{basso2022},
without driving-time rules, and recent work locates chargers for long-haul
trucks under working-hour constraints~\cite{vidotteplaza2026}. The problem
studied here is defined in the companion paper~\cite{pagnier_bakker_journal}.

\section{Problem and Look-Ahead Policy}
\label{sec:problem}

\subsubsection{Route, uncertainty and state.}
A truck drives a fixed route of stops $i\in\{0,\dots,N\}$. The stops are
charging stations $K$, about every 60\,km, laybys $L$, and customers $C$, each
with a service time and possibly a time window $[e_i,l_i]$. Leg $i$ has a
nominal travel time $d_i$ and a realized travel time $d_i\xi_i$, where
$\xi_i=\min\{\xi^{\min}+\eta_i,\xi^{\max}\}$ with
$\eta_i\sim\mathrm{LogNormal}(\mu,\sigma^2)$ independent across legs,
$\xi^{\min}=8/9$, $\xi^{\max}=1.6$, and $(\mu,\sigma)$ such that
$\mathbb{E}[\xi_i]=1$ and the coefficient of variation is 0.15. The travel
time and the energy of a leg are revealed once the leg has been driven. On
arrival at stop $i$, the state $s_i$ consists of the arrival time, the battery
energy and the regulatory counters: driving time since the last break, driving
and on-duty time in the current shift, time since the end of the last daily
rest, the progress of a split break, and the numbers of reduced rests and
extended shifts used.

\subsubsection{Actions and objective.}
At stop $i$, an action $a=(y,b,r)$ specifies whether the truck charges ($y=1$,
only if $i\in K$), and a break $b\in\{b_{45},b_{15},b_{30}\}$, a daily rest
$r\in\{r_{11},r_9\}$, or neither; a charging action is completed by a charging
duration. Driving between breaks is limited to 4.5\,h, driving within a shift
to 9\,h (10\,h twice per week), and a daily rest of 11\,h (9\,h at most three
times) must start within 13\,h (15\,h for a reduced rest) of the end of the
previous one~\cite{eu561,eu2002_15}. We denote by $\mathcal{A}(s)$ the actions
that are structurally legal in state $s$. The battery of 500\,kWh must stay
above 20\% of its capacity, and a station charges at its rated power up to
80\% and more slowly above. A route on which a limit is exceeded is
infeasible. The cost of a route under policy $\pi$ is
\begin{equation}
z(\pi)=t_N+\beta\sum_{i\in C}\mathbf{1}\{t_i\notin[e_i,l_i]\},
\label{eq:obj}
\end{equation}
the arrival time at the destination plus $\beta=0.5$\,h per missed time window,
and the goal is to minimize $\mathbb{E}[z(\pi)]$.

\subsubsection{Look-ahead policy and hindsight oracle.}
Let $\Omega_i$ be 25 travel-time scenarios for the legs within 24\,h of stop
$i$, and $z^H(s_i,a,\omega)$ the optimal value of the MIP
of~\cite{pagnier_bakker_journal} over these legs with the first action fixed
to $a$ and travel times $\omega$. The LA selects
\begin{equation}
a_i^{\mathrm{LA}}\in\arg\min_{a\in\mathcal{A}(s_i)}\ \bar c(s_i,a),\qquad
\bar c(s_i,a)=\frac{1}{|\Omega_i|}\sum_{\omega\in\Omega_i}z^H(s_i,a,\omega),
\label{eq:la}
\end{equation}
where infeasible scenarios are penalized, and executes the first stop of a
nominal re-solve of the MIP. The hindsight \emph{oracle} solves the MIP over
the whole route with the realized travel times; we report the gap
$(z(\pi)-z^\ast)/z^\ast$ to its objective $z^\ast$.

\section{Learned and Hybrid Policies}
\label{sec:method}

\subsubsection{Learning from the look-ahead's action costs.}
To compute~(\ref{eq:la}), the LA evaluates $\bar c(s,a)$ for every
$a\in\mathcal{A}(s)$, and we record all these values rather than its choice
only. Let $\mathcal{A}^F(s)$ be the actions that are feasible in all scenarios.
For $a\in\mathcal{A}^F(s)$, the regression target is the regret
$\rho(s,a)=\bar c(s,a)-\min_{a'\in\mathcal{A}^F(s)}\bar c(s,a')$, and
$\mathbf{1}\{a\in\mathcal{A}^F(s)\}$ is the target of a feasibility model. A
state--action pair is described by 95 features, computed at nominal travel
times: the slack to each regulatory limit, the energy margin to the next
stations, the next three stations and two customers, a summary of the LA's
24-hour window, and the action. We train three models: the regret $\hat c$,
the feasibility $\hat p$, and the charging duration. With gradient-boosted
trees (LightGBM~\cite{ke2017}), each model is a booster, and the regret model
minimizes a Huber loss. With a neural network (PyTorch~\cite{paszke2019}),
each model is a multilayer perceptron, and the regret model minimizes
\begin{equation}
\sum_{s}\Big[\sum_{a\in\mathcal{A}^F(s)}
\ell_\delta\big(\hat c(s,a)-\rho(s,a)\big)
+\lambda\,\mathrm{CE}\big(\mathrm{sm}(-\rho(s,\cdot)/\tau),\,
\mathrm{sm}(-\hat c(s,\cdot)/\tau)\big)\Big],
\label{eq:loss}
\end{equation}
with the Huber loss $\ell_\delta$, the softmax $\mathrm{sm}$ over
$\mathcal{A}^F(s)$, the cross-entropy CE, $\delta=1$\,h, $\tau=0.25$\,h and
$\lambda=1$. The second term penalizes errors in the order of the actions of a
state.

\subsubsection{Safety filter.}
The learned policy selects
$\hat\pi(s)\in\arg\min\{\hat c(s,a): a\in\hat{\mathcal{A}}(s),\
\hat p(s,a)\ge0.5\}$. The set $\hat{\mathcal{A}}(s)\subseteq\mathcal{A}(s)$ is
computed by rules and is never learned: it removes an action if a regulatory
counter would exceed its limit when the next leg is driven at the 0.99
quantile of $\xi$, if the stop would not fit before the next daily rest, or if
the regulation makes another action mandatory. The charging duration is
clipped between the energy needed to reach the next station at the 0.99
quantile and a full battery.

\subsubsection{A gate for calls to the look-ahead.}
The hybrid policy follows $\hat\pi$ and calls the LA at selected stops. Calling
is worth the regret $v(s)=\rho(s,\hat\pi(s))$ of the learned action, which is
known for every state visited by the LA. To obtain representative values on
the training routes, we split them into five folds by route and compute
$\hat\pi(s)$ with a model trained on the other four folds. A gradient-boosted
regressor $g$ predicts $\log(1+v(s))$ from signals that are available online:
the difference between the two best predicted costs, the difference between
the best rest and the best action without rest, the entropy of
$\mathrm{sm}(-\hat c(s,\cdot)/\tau)$, the predicted feasibility, whether a
second learned policy chooses a different action and how much worse the first
one rates that action, and the state features. For a budget of $b\%$ of the
decisions, the LA is called when $g(s)$ exceeds the corresponding quantile on
the validation decisions.

\subsubsection{Rollout check.}
The LA's choice~(\ref{eq:la}) assumes that the LA also takes the subsequent
decisions. In the hybrid policy, $\hat\pi$ does, and an action that is good for
the LA's plan, such as an early daily rest, can be bad when the learned policy
continues. When the LA returns an action $a^{\mathrm{LA}}$ that differs from
$\hat\pi(s)$, we therefore sample $K$ travel-time futures from the LA's
scenario generator, apply either action, and let $\hat\pi$ drive each copy to
the destination. With $\Delta_k$ the cost of the route after $\hat\pi(s)$ minus
the cost after $a^{\mathrm{LA}}$ in future $k$, the action of the LA is
accepted only if $\bar\Delta-z\,\mathrm{se}(\Delta)>0$. Each future is shared
by both actions, so that the comparison is paired. The executed action of the
LA must also belong to $\hat{\mathcal{A}}(s)$, and its charge is raised to the
safe minimum if needed. Algorithm~\ref{alg:hybrid} summarizes one decision.

\begin{algorithm}[t]
\caption{One decision of the hybrid policy at stop $i$}
\label{alg:hybrid}
\begin{algorithmic}[1]
\Require state $s$, learned policy $\hat\pi$, gate $g$ with threshold $\theta_b$, futures $K$, $z$
\State $a\gets\hat\pi(s)$ \Comment{argmin of $\hat c$ over the filtered actions}
\If{$|\hat{\mathcal{A}}(s)|\ge2$ \textbf{and} $g(s)\ge\theta_b$}
  \State $a^{\mathrm{LA}}\gets$ action executed by the LA in state $s$
  \If{$a^{\mathrm{LA}}\in\hat{\mathcal{A}}(s)$ \textbf{and} $a^{\mathrm{LA}}\neq a$}
    \State sample futures $\omega_1,\dots,\omega_K$; for each $k$, drive $\hat\pi$ to the destination after $a$ and after $a^{\mathrm{LA}}$
    \State $\Delta_k\gets$ cost after $a$ minus cost after $a^{\mathrm{LA}}$ under $\omega_k$
    \If{$\bar\Delta-z\,\mathrm{se}(\Delta)>0$} $a\gets a^{\mathrm{LA}}$ \EndIf
  \EndIf
\EndIf
\State \Return $a$
\end{algorithmic}
\end{algorithm}

\section{Computational Experiments}
\label{sec:exp}

Our experiments consist of three parts. First, we compare the learned policies
with the LA on routes of the kind they were trained on. Second, we evaluate the
hybrid policy on the same routes, together with ablations of the gate and of
the rollout check. Third, we evaluate the hybrid policy on routes whose
charging stations have different powers.

\subsection{Setup}
\label{sec:setup}

\subsubsection{Instances.}
Routes are 800--1{,}200, 1{,}500--2{,}500 or 3{,}000--4{,}000\,km long and
serve 1--3, 4--6 or 7--15 customers with no, tight, medium or large time
windows. Each of the resulting 36 families has 25 routes with one travel-time
realization each, and all stations have 350\,kW. Routes with seeds 1--19 are
used for training (639 routes solved by the LA, 56{,}104 decisions), seeds
20--21 for validation (66 routes), and seeds 22--25 for testing (125 of the
144 routes have an LA run). For the third part, the power of each station is
drawn from $\{150,350,700,1000\}$\,kW, either alone or together with the
distance between stations; the LA was run on 238 such training routes and on
64 test routes, and 32 routes serve for validation.

\subsubsection{Measures and protocol.}
We report the gap to the oracle with the time-window penalty
of~(\ref{eq:obj}) and its difference with the LA in percentage points (pp),
paired by route; $\pm$ is one standard error over routes, and the results of
several training seeds are averaged per route first. The oracle is solved to a
relative optimality gap of 0.5\% on 114 of the 125 test routes. All design
choices were made on the validation routes; each configuration reported on the
test routes was run once. The LA uses Gurobi~12 through Pyomo with eight
parallel workers; the learned policies run on a CPU.

\subsection{Learned Policies}
\label{sec:learned}

% numbers: ML/scratch/both_metrics.py machinery on eval_*_g99srNm_test.json
% (3 training seeds), nsN for the runs without filter; compare_oracle.py
\begin{table}[t]
\caption{Learned policies on the 125 test routes (three training seeds each):
difference with the LA in percentage points of the gap to the oracle, with and
without the time-window penalty, missed time windows over all routes (737
customer visits), and runs that violated the regulation. The LA is 3.05\%
above the oracle (2.77\% without penalty). The Greedy rule is the heuristic
of~\cite{pagnier_bakker_journal}.}
\label{tab:learned}
\centering
\begin{tabular}{lrrrr}
\toprule
Policy & vs.\ LA & without penalty & missed windows & violations\\
\midrule
LA & -- & -- & 57 & 0/125\\
Trees & $+0.68\pm0.43$ & $+0.23\pm0.41$ & 116 & 0/375\\
Neural network & $+0.50\pm0.37$ & $+0.10\pm0.35$ & 111 & 0/375\\
Greedy & $+3.95\pm0.46$ & $+2.81\pm0.40$ & 174 & 5/125\\
\bottomrule
\end{tabular}
\end{table}

Table~\ref{tab:learned} shows that both learned policies are within 0.7\,pp of
the LA, and within 0.25\,pp when the time-window penalty is ignored: they miss
twice as many time windows as the LA. A decision takes 1 to 5\,ms, against more
than a minute for the LA. No run violates the regulation, which is due to the
safety filter rather than to the models: without the filter, 18 of the 375
runs of the trees and 38 of the 375 runs of the network violate a limit.

Using the oracle as teacher, i.e., training the same models to reproduce the
oracle's action at every stop of the stored hindsight-optimal schedules,
yields worse policies, by $+2.05\pm0.48$\,pp for the trees and
$+6.42\pm0.66$\,pp for the network. These policies are as good as the LA on
most routes but take one more daily rest than the LA on 13\% (trees) and 32\%
(network) of the runs, against 3\% and 2\% for the same models trained on the
LA's costs. The oracle plans the last shifts close to the limits
because it knows the travel times, and a policy that reproduces these choices
without this knowledge is regularly forced into an additional rest.

\subsection{Calling the Look-Ahead}
\label{sec:hybrid}

We study the hybrid policy with the trees, which leave more room to the LA than
the network on the validation routes. On the validation decisions, the gate
recovers 65\% of the regret of the learned policy by calling the LA at 5\% of
the decisions, and 87\% at 10\%, whereas calling where the best rest and the
best non-rest action are closest recovers 26\% and 30\%, respectively. In
closed loop, however, the gate alone helped less than these numbers suggest.
On 2 of the 66 validation routes, a call to the LA led to an additional daily
rest, adding 16 and 30\,pp to these routes. On the second route, the LA chose
an early daily rest; the learned policy, which took the subsequent decisions,
then reached the 15-hour limit and had to rest a second time. The rollout check
was introduced for this reason; with $K=32$ and $z=1$, it costs about 30\,s of
simulation per route.

% numbers: ML/code/gate_report.py --set test, hybrid_base_*L_ms_N.jsonl
\begin{table}[t]
\caption{Hybrid policies on the 125 test routes (trees, one training seed), run
once with the configuration chosen on the validation routes: calls to the LA
per route (out of 95.5 decisions), difference with the LA and with the learned
policy alone in percentage points, routes with more daily rests than the LA,
and violations. All hybrid variants filter the LA's actions.}
\label{tab:hybrid}
\centering
\begin{tabular}{lrrrrr}
\toprule
Policy & calls & vs.\ LA & vs.\ alone & extra rests & violations\\
\midrule
Learned policy alone & 0 & $+0.73\pm0.44$ & -- & 3 & 0\\
Gate, without check & 9.1 & $+0.09\pm0.35$ & $-0.64\pm0.30$ & 2 & 0\\
Rest margin with check & 9.4 & $+0.00\pm0.37$ & $-0.73\pm0.26$ & 2 & 0\\
Gate with check & 8.1 & $-0.36\pm0.33$ & $-1.09\pm0.30$ & 2 & 0\\
\bottomrule
\end{tabular}
\end{table}

Table~\ref{tab:hybrid} reports the test results. The hybrid policy with gate
and rollout check is at the level of the LA ($-0.36\pm0.33$\,pp) and improves
on the learned policy by $1.09\pm0.30$\,pp, while calling the LA at 8.5\% of
the decisions; including the roll-outs, this is about 10\% of the LA's
computing time. At a similar number of calls, the gate is better than the rest
margin ($-0.37\pm0.18$\,pp, paired) and the rollout check improves the gate
($-0.45\pm0.11$\,pp); the check rejects 33\% of the LA's actions. The filter
rejected 23 actions of the LA that it judged unsafe, and no run violates the
regulation. On the validation routes, the same configuration was ahead of the
LA ($-0.67\pm0.28$\,pp); the test routes do not confirm an advantage over the
LA, and the gain over the learned policy is concentrated on short and medium
routes.

\subsection{Charging Stations of Different Powers}
\label{sec:mixed}

When the stations of a route have different powers, the LA loses nothing
compared with the same route with equal powers, because it charges mostly at
the fast stations. Learned policies trained on routes with equal powers do not
learn to choose between stations, and adding the LA's runs on 238 training
routes with varying powers leaves the network $2.84\pm1.04$\,pp behind the LA
on the 61 test routes with a certified oracle ($1.94\pm1.10$ when only the
power varies, $3.76\pm1.79$ when the distance varies as well). For this
learned policy, the second policy of the gate is the same network trained with
another seed, since the LA was not run on the validation routes, the gate is
selected on held-out training routes. On the 30 validation routes, calling the
LA where the rest margin is below 20 minutes improves on the learned policy by
$0.45\pm0.25$\,pp with the rollout check, and is worse without it (one
additional daily rest). \tbd{Gate with check on the validation routes; test
results against the LA in a table like Table~\ref{tab:hybrid}.}

\section{Conclusion}
\label{sec:conclusion}

We studied learned policies for the charging, break and rest decisions of an
electric truck, trained on the action costs of a stochastic look-ahead, and a
hybrid policy that calls the look-ahead only where a learned gate expects a
gain. On routes of the kind they were trained on, the learned policies are
close to the look-ahead and, with a rule-based filter, never violate the
regulation; reproducing the hindsight-optimal schedules instead leads to
additional daily rests. The hybrid policy reaches the look-ahead at a tenth of
its computing time, provided that each action of the look-ahead is checked by
rolling out the learned policy, which will take the subsequent decisions.
\tbd{Mixed-power conclusion.} The main limitations are the synthetic
instances, one training seed for the hybrid policies, and the remaining errors
of the look-ahead itself, in particular additional daily rests near the end of
long routes, which neither the learned nor the hybrid policies remove. A
natural extension is to learn the gate from the rolled-out value of a call
rather than from the look-ahead's own cost, which our first attempt did not
improve, and to train the learned policy with roll-outs of its own.

\begin{credits}
\subsubsection{\ackname}
% TODO: funding
\subsubsection{\discintname}
The authors have no competing interests to declare that are relevant to the
content of this article.
\end{credits}
%
% CPAIOR 2027 call: "Submissions representing work that has been submitted to
% or published in another venue must explicitly state that venue." -> name the
% venue of the companion paper in the text or a footnote (TODO).
\bibliographystyle{splncs04}
\bibliography{references}
\end{document}
